\SourceLeaf{32}
\fancyhead[L]{\small BEOBACHTUNGEN DES ZWEITEN KOMETEN VON 1813}
\renewcommand{\Article}[1]{\par\medskip\Needspace{4\baselineskip}\pdfbookmark[1]{Art. #1}{cometae-art-#1}\centerline{\textbf{#1.}}\nobreak\smallskip}
\pdfbookmark[0]{Beobachtungen des zweiten Kometen von 1813}{cometae}
\begin{center}
{\large BEOBACHTUNGEN}\par\medskip
{\Large DES ZWEITEN KOMETEN VON 1813}\par\medskip
ANGESTELLT AN DER GÖTTINGER STERNWARTE\par\medskip
MIT\par\medskip
EINIGEN BEMERKUNGEN ZUR BERECHNUNG PARABOLISCHER BAHNEN.\par\smallskip
\rule{.09\linewidth}{.3pt}
\end{center}

Am 7. April begann ich selbst, den Kometen an unserer Sternwarte zu beobachten, nachdem ihn mein lieber Kollege, der angesehene \textsc{Harding}, am 3. April dieses Jahres im Sternbild des Poniatowskischen Stiers entdeckt hatte. Mit einem Kreismikrometer, das an einem zehnfüßigen Fernrohr angebracht war, erhielt ich folgende Bestimmungen:
\begin{center}\small\setlength{\tabcolsep}{4pt}
\begin{tabular}{r|r|r|r l}
1813&Mittl. Gött. Zeit&Scheinb. Rektasz.&\multicolumn{2}{c}{Scheinb. Dekl.}\\\hline
Apr. 7&$13^{\mathrm h}\;12^{\mathrm m}\;2^{\mathrm s}$&$271^0\;7'\;19''3$&$5^0\;34'\;36''7$&Nördl.\\
9&$13\;35\;40$&$270\;10\;33.5$&$4\;11\;3.4$&—\\
11&$13\;17\;43$&$269\;1\;19.9$&$2\;33\;0.7$&—\\
14&$13\;7\;36$&$266\;44\;5.5$&$0\;33\;0.8$&Südl.\\
21&$14\;23\;0$&$256\;39\;19.3$&$12\;57\;56.0$&—
\end{tabular}\end{center}

Der angesehene \textsc{Harding} stellte später mit dem Mauerquadranten folgende Beobachtungen an:
\begin{center}\small\setlength{\tabcolsep}{4pt}
\begin{tabular}{r|r|r|r l}
Apr. 21&$15^{\mathrm h}\;7^{\mathrm m}\;21^{\mathrm s}$&$256^0\;34'\;19''6$&$13^0\;2'\;26''5$&Südl.\\
24&$14\;22\;50$&$248\;23\;21$&$21\;45\;2$&\\
25&$14\;4\;21$&$244\;44\;42$&$25\;10\;42$&
\end{tabular}\end{center}

Am 24. und 25. April war der Komet auch mit bloßem Auge sehr deutlich zu sehen. In den folgenden Nächten war der Himmel von Wolken bedeckt, und bald setzte der rasche Abstieg des Kometen nach Süden den Beobachtungen ein Ende.

\SourceLeaf{33}\SourceGathering{4\textsuperscript{*}}
Es scheint überflüssig, hier die parabolischen Elemente zu wiederholen, die ich gleich zu Beginn selbst aus den ersten drei Beobachtungen abgeleitet hatte. Mit ihrer genaueren Ausarbeitung betraute ich einen sehr geübten und erfahrenen Rechner, Dr. \textsc{Gerling}. Ihm verdanken wir die folgenden korrigierten Elemente, die so gut wie möglich an alle Beobachtungen angepasst sind, sowohl an unsere eigenen als auch an die, die der angesehene \textsc{Olbers} übersandt hatte.
\par\medskip
\noindent\begin{minipage}{.975\linewidth}\small
\noindent\makebox[\linewidth][s]{Logarithmus der Periheldistanz\hfill.\hfill.\hfill.\hfill.\hfill.\hfill.\hfill.\hfill.\hfill.\hfill.\hfill$0.0849212$}\par
\noindent Zeit des Periheldurchgangs für den Göttinger Meridian\par
\noindent\makebox[\linewidth][r]{1813 Mai $19.44507$}\par
\noindent\makebox[\linewidth][s]{Länge des Perihels\hfill.\hfill.\hfill.\hfill.\hfill.\hfill.\hfill.\hfill.\hfill.\hfill.\hfill.\hfill.\hfill.\hfill.\hfill$197^0\;43'\;7''7$}\par
\noindent\makebox[\linewidth][s]{Länge des aufsteigenden Knotens\hfill.\hfill.\hfill.\hfill.\hfill,\hfill.\hfill.\hfill.\hfill.\hfill.\hfill.\hfill$42\;40\;15.2$}\par
\noindent\makebox[\linewidth][s]{Bahnneigung\hfill.\hfill.\hfill.\hfill.\hfill.\hfill.\hfill.\hfill.\hfill.\hfill.\hfill.\hfill.\hfill.\hfill.\hfill.\hfill$81\;2\;11.8$}\par
\noindent \textit{Rückläufige} Bewegung.\par
\end{minipage}\par\medskip

\Needspace{11\baselineskip}
Der angesehene \textsc{Olbers} stellte folgende Beobachtungen an:
\begin{center}\small\setlength{\tabcolsep}{4pt}
\begin{tabular}{r|r|r|r l}
1813&Mittl. Bremer Zeit&Scheinb. Rektasz.&\multicolumn{2}{c}{Scheinb. Dekl.}\\\hline
Apr. 14&$13^{\mathrm h}\;31^{\mathrm m}\;4^{\mathrm s}$&$266^0\;42'\;51''2$&$0^0\;34'\;22''8$&Südl.\\
15&$12\;14\;29$&$265\;48\;47.9$&$1\;46\;4.5$&\\
19&$11\;38\;0$&$260\;40\;39.1$&$8\;15\;23.7$&\\
21&$12\;0\;35$&$256\;51\;59.3$&$12\;42\;54.3$&\\
24&$11\;58\;38$&$248\;43\;57.7$&$21\;25\;9.8$&\\
25&$11\;41\;30$&$245\;8\;18.0$&$24\;49\;2.4$&\\
&$12\;5\;38$&$245\;4\;3.0$&$24\;54\;16.4$&
\end{tabular}\end{center}

Hinzu kam eine Beobachtung des angesehenen \textsc{Bouvard} an der Kaiserlichen Sternwarte in Paris:
\begin{center}\small\setlength{\tabcolsep}{4pt}
\begin{tabular}{r|r|r|r l}
1813&Mittl. Pariser Zeit&Scheinb. Rektasz.&\multicolumn{2}{c}{Scheinb. Dekl.}\\\hline
Apr. 13&$16^{\mathrm h}\;22^{\mathrm m}\;2^{\mathrm s}$&$267^0\;27'\;18''$&$0^0\;24'\;46''$&Nördl.
\end{tabular}\end{center}

Die Rechnungen des angesehenen \textsc{Gerling} zeigen, wie gut die oben angegebenen Elemente mit den einzelnen Beobachtungen übereinstimmen:
\SourceLeaf{34}
\begin{center}\small\setlength{\tabcolsep}{6pt}
\begin{tabular}{r|r|r|l}
&\multicolumn{2}{c|}{Differenz}&\\
&Rektasz.&Dekl.&Beobachter\\\hline
Apr. 7&$+3''8$&$+8''5$&Gauss\\
9&$+2.0$&$+34.3$&Gauss\\
11&$-5.3$&$-17.7$&Gauss\\
13&$-1.6$&$-0.4$&Bouvard\\
14&$-7.4$&$-28.6$&Gauss\\
&$+2.7$&$-8.4$&Olbers\\
15&$-0.9$&$+28.7$&Olbers\\
19&$-25.1$&$+103.9$&Olbers\\
21&$-56.6$&$-59.2$&Olbers\\
&$-30.1$&$-5.2$&Gauss\\
&$-22.8$&$-24.1$&Harding\\
24&$-45.2$&$-41.4$&Olbers\\
&$+0.4$&$-11.6$&Harding\\
25&$-23.3$&$-67.8$&Olbers\\
&$-9.4$&$-27.4$&Olbers\\
&$+9.7$&$+1.4$&Harding
\end{tabular}\end{center}

Bei diesem Vergleich wurden sowohl die Aberration als auch die Parallaxe gebührend berücksichtigt.
\par\nobreak\smallskip
\noindent\makebox[\linewidth][c]{\rule{.225\linewidth}{.3pt}}\par\medskip

Ich möchte einige Rechenabkürzungen hinzufügen, die ich bei der ersten Bestimmung einer parabolischen Bahn nach der Methode des angesehenen \textsc{Olbers} oft mit Vorteil benutzt habe. Mit ihnen lässt sich diese ohnehin schon rasche Methode, wie es scheint, weiter abkürzen oder besser an die numerische Rechnung anpassen. Sie betreffen die Berechnung der Radiusvektoren und vor allem die Sehne zwischen dem ersten und dem letzten Ort. Der angesehene \textsc{Olbers} verwendet dazu Ausdrücke der Form $\BareRoot{(f+g\rho+h\rho\rho)}$. Die Koeffizienten $f$, $g$, $h$ bestimmt er durch recht einfache Rechenschritte. In den meisten Fällen erreicht man damit jedoch nur dann ausreichende Genauigkeit, wenn man sie mit den größeren Logarithmentafeln ausführt, die auf sieben oder wenigstens sechs Dezimalstellen berechnet sind. Ich habe dafür andere Ausdrücke eingesetzt, die sich an sich etwas leichter berechnen lassen und außerdem den Vorteil haben, dass für alle Rechenschritte fünf Dezimalstellen und damit die kleineren Logarithmentafeln genügen. Das Wesentliche ist Folgendes. Es seien
\SourceLeaf{35}
\par\noindent\hspace*{.06\linewidth}%
\begin{tabular}[t]{@{}l@{\enspace}p{.70\linewidth}@{}}
  $\odot$, $\odot'$, $\odot''$ & die Längen der Sonne bei der ersten, zweiten und dritten Beobachtung\\
  $R$, $R'$, $R''$ & die Entfernungen der Sonne von der Erde\\
  $\alpha$, $\alpha'$, $\alpha''$ & die geozentrischen Längen und\\
  $\ObservedBeta$, $\ObservedBeta'$, $\ObservedBeta''$ & die geozentrischen Breiten des Kometen\\
  $r$, $r'$, $r''$ & seine Entfernungen von der Sonne\\
  $\rho$, $\rho'$, $\rho''$ & seine verkürzten Entfernungen von der Erde\\
  $t$, $t'$, $t''$ & die Beobachtungszeiten\\
  \multicolumn{2}{@{}l@{}}{$k$\quad die Sehne zwischen dem ersten und dem dritten Ort des Kometen}\\
  \multicolumn{2}{@{}l@{}}{$M=\frac{\rho''}{\rho}$}
\end{tabular}\par\medskip

Damit sieht man leicht, dass gilt
\[
\begin{aligned}{}
{[1]}\quad r&=\BareRoot{[(\rho\cos\alpha-R\cos\odot)^2}\\
&\qquad +(\rho\sin\alpha-R\sin\odot)^2+\rho\rho\operatorname{tang}\ObservedBeta^2]\\
{[2]}\quad r''&=\BareRoot{[(M\rho\cos\alpha''-R''\cos\odot'')^2}\\
&\qquad +(M\rho\sin\alpha''-R''\sin\odot'')^2\\
&\qquad +MM\rho\rho\operatorname{tang}{\ObservedBeta''}^2]\\
{[3]}\quad k&=\BareRoot{[(M\rho\cos\alpha''-\rho\cos\alpha-R''\cos\odot''+R\cos\odot)^2}\\
&\qquad +(M\rho\sin\alpha''-\rho\sin\alpha-R''\sin\odot''+R\sin\odot)^2\\
&\qquad +(M\rho\operatorname{tang}\ObservedBeta''-\rho\operatorname{tang}\ObservedBeta)^2]
\end{aligned}
\]
Die Gleichungen 1 und 2 nehmen folgende Form an
\[
\begin{aligned}
r&=\BareRoot{\left(\frac{\rho\rho}{\cos\ObservedBeta^2}-2\rho R\cos(\alpha-\odot'')+RR\right)}\\
r''&=\BareRoot{\left(\frac{MM\rho\rho}{\cos{\ObservedBeta''}^2}-2M\rho R''\cos(\alpha''-\odot'')+R''R''\right)}
\end{aligned}
\]
Setzen wir also
\[
\begin{aligned}
\cos\ObservedBeta\cos(\alpha-\odot)&=\cos\psi, & R\sin\psi&=B\\
\cos\ObservedBeta''\cos(\alpha''-\odot'')&=\cos\psi'', & R''\sin\psi''&=B''
\end{aligned}
\]
so erhalten wir
\[
\begin{aligned}
r&=\BareRoot{\left[\left(\frac{\rho}{\cos\ObservedBeta}-R\cos\psi\right)^2+BB\right]}\\
r''&=\BareRoot{\left[\left(\frac{M\rho}{\cos\ObservedBeta''}-R''\cos\psi''\right)^2+B''B''\right]}
\end{aligned}
\]

Bestimmen wir ferner die fünf Hilfsgrößen $g$, $G$, $h$, $H$, $\zeta$ so, dass gilt
\SourceLeaf{36}
\[
\begin{aligned}
R''\cos\odot''-R\cos\odot&=g\cos G\\
R''\sin\odot''-R\sin\odot&=g\sin G\\
M\cos\alpha''-\cos\alpha&=h\cos\zeta\cos H\\
M\sin\alpha''-\sin\alpha&=h\cos\zeta\sin H\\
M\operatorname{tang}\ObservedBeta''-\operatorname{tang}\ObservedBeta&=h\sin\zeta
\end{aligned}
\]
so geht Formel 3 in die folgende über
\[
\begin{aligned}
k&=\BareRoot{[(\rho h\cos\zeta\cos H-g\cos G)^2}\\
&\qquad +(\rho h\cos\zeta\sin H-g\sin G)^2+\rho\rho hh\sin\zeta^2]\\
&=\BareRoot{(\rho\rho hh-2\rho hg\cos\zeta\cos(G-H)+gg)}
\end{aligned}
\]
Setzen wir also
\[
\cos\zeta\cos(G-H)=\cos\varphi,\quad g\sin\varphi=A
\]
so gilt
\[
k=\BareRoot{[(\rho h-g\cos\varphi)^2+AA]}
\]
oder, wenn wir zusätzlich $\rho h-g\cos\varphi=u$ setzen,
\[
k=\BareRoot{(uu+AA)}
\]

Wir glauben, dass es vielen Lesern willkommen sein wird, wenn wir hier sämtliche zu diesen Umformungen gehörenden Rechenschritte nochmals in der richtigen Reihenfolge darstellen und auch alle übrigen Schritte hinzufügen. So findet sich hier alles beisammen, was zur ersten Berechnung einer parabolischen Bahn erforderlich ist. Zugleich erläutern wir diese Vorschriften mit Zahlen zu unserem Kometen. Wir wählen dazu unsere Beobachtungen vom 7., 14. und 21. April, deren Reduktion die folgenden \textit{Daten} liefert:
\[
\begin{array}{rcrrr@{\qquad}rcrrr}
t&=&\multicolumn{3}{l}{7.55002}&&&&&\\
t'&=&\multicolumn{3}{l}{14.54694}&&&&&\\
t''&=&\multicolumn{3}{l}{21.59931}&&&&&\\
\alpha&=&271^0&16'&38''&\ObservedBeta&=&+29^0&2'&0''\\
\alpha'&=&266&27&22&\ObservedBeta'&=&+22&52&18\\
\alpha''&=&256&48&8&\ObservedBeta''&=&+9&53&12\\
\odot&=&17&47&41&\log R&=&\multicolumn{3}{l}{0.00091}\\
\odot'&=&24&38&45&\log R'&=&\multicolumn{3}{l}{0.00175}\\
\odot''&=&31&31&25&\log R''&=&\multicolumn{3}{l}{0.00260}
\end{array}
\]

\SourceLeaf{37}
I. Der \textit{erste} Rechenschritt besteht darin, einen Näherungswert für $M$ zu bestimmen. Ihn liefert die Formel
\[
M=\frac{t''-t'}{t'-t}.\frac{\operatorname{tang}\ObservedBeta'\sin(\alpha-\odot')-\operatorname{tang}\ObservedBeta\sin(\alpha'-\odot')}{\operatorname{tang}\ObservedBeta''\sin(\alpha'-\odot')-\operatorname{tang}\ObservedBeta'\sin(\alpha''-\odot')}
\]
In unserem Beispiel ergibt sich $\log M=9.75799$.

II. Nun sind die Größen $g$, $G$, $h$, $H$, $\zeta$ nach den folgenden Formeln zu bestimmen. Sie sind offensichtlich den oben angegebenen gleichwertig und für die Rechnung bequemer.
\[
\begin{aligned}
R''\cos(\odot''-\odot)-R&=g\cos(G-\odot)\\
R''\sin(\odot''-\odot)&=g\sin(G-\odot)\\
M-\cos(\alpha''-\alpha)&=h\cos\zeta\cos(H-\alpha'')\\
\sin(\alpha''-\alpha)&=h\cos\zeta\sin(H-\alpha'')\\
M\operatorname{tang}\ObservedBeta''-\operatorname{tang}\ObservedBeta&=h\sin\zeta
\end{aligned}
\]

In unserem Beispiel ergibt sich
\[
\begin{array}{rcrrr}
G&=&113^0&43'&57''\\
\log g&=&\multicolumn{3}{l}{9.38029}\\
H&=&109^0&5'&49''\\
\zeta&=&44&13&9\\
\log h&=&\multicolumn{3}{l}{9.81477}
\end{array}
\]

III. Danach setzen wir
\[
\begin{aligned}
\cos\zeta\cos(G-H)&=\cos\varphi\\
\cos\ObservedBeta\cos(\alpha-\odot)&=\cos\psi\\
\cos\ObservedBeta''\cos(\alpha''-\odot'')&=\cos\psi''\\[5pt]
g\sin\varphi&=A\\
R\sin\psi&=B\\
R''\sin\psi''&=B''
\end{aligned}
\]
Wenn sich bei dieser Rechnung die Kosinuswerte der Winkel $\varphi$, $\psi$, $\psi''$ nur wenig von eins unterscheiden, empfiehlt es sich, sechs oder sogar sieben Stellen zu verwenden. Die Winkel $\varphi$, $\psi$, $\psi''$ selbst braucht man nicht in Grad, Minuten und Sekunden anzugeben, da es genügt, von den Logarithmen der Kosinuswerte unmittelbar zu denen der Sinuswerte überzugehen.

\SourceLeaf{38}
In unserem Beispiel ergibt sich
\[
\begin{aligned}
\log A&=9.22527\\
\log B&=9.98706\\
\log B''&=9.86038
\end{aligned}
\]

IV. Schließlich setze man
\[
\begin{aligned}
h\cos\ObservedBeta&=b\\
\frac{h\cos\ObservedBeta''}{M}&=b''\\
g\cos\varphi-bR\cos\psi&=c\\
g\cos\varphi-b''R''\cos\psi''&=c''
\end{aligned}
\]
In unserem Beispiel
\[
\begin{aligned}
\log b&=9.75645\\
\log b''&=0.05028\\
c&=+0.31365\\
c''&=+0.95443
\end{aligned}
\]

V. Nach diesen Vorbereitungen hängen die Radiusvektoren $r$, $r''$ und die Sehne $k$ wie folgt von der Unbekannten $u$ ab
\[
\begin{aligned}
r&=\BareRoot{\left[\left(\frac{u+c}{b}\right)^2+BB\right]}\\
r''&=\BareRoot{\left[\left(\frac{u+c''}{b''}\right)^2+B''B''\right]}\\
k&=\BareRoot{(uu+AA)}
\end{aligned}
\]
Die Unbekannte ist durch Probieren so zu bestimmen, dass die folgende Gleichung erfüllt ist
\[
(r+r''+k)^{\frac32}-(r+r''-k)^{\frac32}=\frac{t''-t}{m}
\]
Dabei bezeichnet $m$ eine Zeit von $9.6887401$ Tagen, und es gilt
\[
\log m=0.9862673
\]
Der Größe $(r+r''-k)^{\frac32}$ müsste das Zeichen $+$ vorangestellt werden, wenn die heliozentrische Bewegung des Kometen in der Zeit $t''-t$ einen Winkel von $180^0$ überschritte. Dieser Fall kann unter den Voraussetzungen, auf denen diese erste Bahnbestimmung beruht, nicht eintreten. Es braucht wohl kaum erwähnt zu werden, dass man zur numerischen Berechnung von $r$ einen Hilfswinkel $\theta$ einführen sollte, für den gilt
\SourceLeaf{39}\SourceGathering{VI.~5}
\[
\frac{bB}{u+c}=\operatorname{tang}\theta
\]
Daraus folgt $r=\frac{B}{\sin\theta}$, und ebenso verfährt man bei $r''$ und $k$. Jeder wird außerdem leicht erkennen, dass sich bei all diesen Rechenschritten unsere Tafel zur unmittelbaren Bestimmung der Logarithmen von Summen und Differenzen sehr bequem verwenden lässt.

In unserem Beispiel gilt $\log\frac{t''-t}{m}=0.16139$, und nach wenigen Versuchen ergibt sich
\[
u=0.24388
\]

VI. Hat man $u$ gefunden, so gilt
\[
\rho=\frac{u+g\cos\varphi}{h},\quad\rho''=M\rho
\]
\begin{center}(in unserem Beispiel $\log\rho=9.80364$, $\log\rho''=9.56163$).\end{center}

Die übrigen Rechenschritte sind hinreichend bekannt. Damit aber hier alles beisammen ist, schien es zweckmäßig, auch die übrigen Formeln anzufügen, die wir gewöhnlich verwenden. Es seien
\begin{description}
\setlength{\itemsep}{1.2pt}
\setlength{\parsep}{0pt}
\item[$\lambda$, $\lambda''$] die heliozentrischen Längen des Kometen bei der ersten und dritten Beobachtung
\item[$\ApproxBeta$, $\ApproxBeta''$] die heliozentrischen Breiten
\item[$\upsilon$, $\upsilon''$] die Längen in der Bahn
\item[$\AscendingNode$] die Länge des aufsteigenden Knotens
\item[$i$] die Bahnneigung, die zwischen $0$ und $90^0$ anzusetzen ist, wenn wir wie üblich rechtläufige und rückläufige Bewegung unterscheiden
\item[$\Pi$] die Länge des Perihels
\item[$T$] die Zeit des Periheldurchgangs
\item[$q$] die Periheldistanz.
\end{description}

VII. Die heliozentrischen Orte findet man mit den Formeln
\[
\begin{aligned}
\rho\cos(\alpha-\odot)-R&=r\cos\ApproxBeta\cos(\lambda-\odot)\\
\rho\sin(\alpha-\odot)&=r\cos\ApproxBeta\sin(\lambda-\odot)\\
\rho\operatorname{tang}\ObservedBeta&=r\sin\ApproxBeta\\
\rho''\cos(\alpha''-\odot'')-R''&=r''\cos\ApproxBeta''\cos(\lambda''-\odot'')\\
\rho''\sin(\alpha''-\odot'')&=r''\cos\ApproxBeta''\sin(\lambda''-\odot'')\\
\rho''\operatorname{tang}\ObservedBeta''&=r''\sin\ApproxBeta''
\end{aligned}
\]

\SourceLeaf{40}
Die Übereinstimmung der hier erhaltenen Werte für die Radiusvektoren $r$, $r''$ mit den zuvor aus $u$ abgeleiteten Werten dient zur Kontrolle der Rechnung. Die Bewegung ist rechtläufig oder rückläufig, je nachdem, ob $\lambda''$ größer oder kleiner als $\lambda$ ist.

In unserem Beispiel finden wir
\[
\begin{aligned}
\lambda&=225^0\;4'\;22'',&\ApproxBeta&=+14^0\;51'\;39'',&\log r&=0.13896\\
\lambda''&=223\;6\;55,&\ApproxBeta''&=+2\;49\;28,&\log r''&=0.11068
\end{aligned}
\]

Die Bewegung des Kometen ist also rückläufig.

VIII. Zur Bestimmung der Knotenlänge und der Bahnneigung verwendet man die Formeln:
\[
\begin{aligned}
\pm\operatorname{tang}\ApproxBeta&=\operatorname{tang}i.\sin(\lambda-\AscendingNode)\\
\pm\frac{\operatorname{tang}\ApproxBeta''-\operatorname{tang}\ApproxBeta.\cos(\lambda''-\lambda)}{\sin(\lambda''-\lambda)}&=\operatorname{tang}i.\cos(\lambda-\AscendingNode)
\end{aligned}
\]
wobei die oberen Zeichen für rechtläufige, die unteren für rückläufige Bewegung gelten. Die Längen in der Bahn ergeben sich dann aus den Formeln
\[
\begin{aligned}
\frac{\operatorname{tang}(\lambda-\AscendingNode)}{\cos i\cdot}&=\operatorname{tang}(\upsilon-\AscendingNode)\\
\frac{\operatorname{tang}(\lambda''-\AscendingNode)}{\cos i}&=\operatorname{tang}(\upsilon''-\AscendingNode)
\end{aligned}
\]
wobei $\upsilon-\AscendingNode$, $\upsilon''-\AscendingNode$ jeweils in denselben Quadranten zu nehmen sind wie $\lambda-\AscendingNode$, $\lambda''-\AscendingNode$.

Für unseren Kometen finden wir
\[
\begin{array}{rcrrr}
\AscendingNode&=&42^0&40'&8''\\
i&=&81&1&3\\
\upsilon&=&237&43&7\\
\upsilon''&=&225&31&32
\end{array}
\]

IX. Die Länge des Perihels und die Periheldistanz ergeben sich aus den Formeln
\[
\begin{gathered}
\frac{1}{\BareRoot r}=\frac{1}{\BareRoot q}.\cos\tfrac12(\upsilon-\Pi)\\
\frac{\operatorname{cotg}\tfrac12(\upsilon''-\upsilon)}{\BareRoot r}-\frac{1}{\sin\tfrac12(\upsilon''-\upsilon).\BareRoot{r''}}=\frac{1}{\BareRoot q}.\sin\tfrac12(\upsilon-\Pi)
\end{gathered}
\]
Für unseren Kometen gilt $\Pi=197^0\;37'\;51''$, $\log q=0.08469$.

\SourceLeaf{41}\SourceGathering{5\textsuperscript{*}}
X. Schließlich entnehme man der Tafel von \textsc{Barker} die mittleren Bewegungen, die den wahren Anomalien $\upsilon-\Pi$, $\upsilon''-\Pi$ oder $\Pi-\upsilon$, $\Pi-\upsilon''$ entsprechen, und bezeichne sie mit $M$, $M''$. Dann gilt
\[
T=t\mp Mnq^{\frac32}=t''\mp M''nq^{\frac32}
\]
Die oberen Zeichen gelten bei rechtläufiger Bewegung, wenn $\upsilon>\Pi$, $\upsilon''>\Pi$, oder bei rückläufiger Bewegung, wenn $\upsilon<\Pi$, $\upsilon''<\Pi$. In den entgegengesetzten Fällen gelten die unteren Zeichen. Die Größe $n$ ist konstant, mit dem Logarithmus $=0.0398723$. Die Übereinstimmung der beiden Werte von $T$ liefert eine zweite Kontrolle der Rechnung.

In unserem Beispiel finden wir
\[
\begin{aligned}
T&=49.518\\
T&=49.517
\end{aligned}
\]
sodass der $19.5175$. Mai als Zeitpunkt des Periheldurchgangs angesetzt werden kann.

Berechnet man aus diesen Elementen den geozentrischen Ort zur Zeit der mittleren Beobachtung, so erhält man die Länge $266^0\;27'\;15''$ und die nördliche Breite $22^0\;52'\;18''$. Die Länge weicht um $7''$ von der beobachteten ab, die Breite stimmt genau überein.
\par\nobreak\medskip
\centerline{\rule{.25\linewidth}{.3pt}}
